\section{嘉道危机证据链事件锚点}
以下锚点为嘉道时期战争、河工、边贸、禁烟和条约补入独立证据节点，防止行政文本取代地方社会与跨境材料。
\subsection{1790—1799}
\eventanchor{EQ-1796-JIAQING-ACCESSION-DOCUMENTS}{围绕「嘉庆帝正式即位，白莲教战…}嘉庆元年（1796年），围绕「嘉庆帝正式即位，白莲教战争与财政压力成为初政核心问题」的上谕、奏折与部议在中央—地方行政链中流转。核心取证：《清仁宗实录》嘉庆元年相关条；宫中朱批奏折、内阁大库题本与《清史稿》相关志传。这一锚点记录的是文书链：它能证明呈报、上谕、部议或照会进入机构流程，不能由此推出实际执行。来源保留：此行仅确认相关文书链和可追查的行政动作；不得由其直接推断政策有效、地方服从、民众态度或单一因果。
\eventanchor{EQ-1797-WHITE-LOTUS-WAR-DOCUMENTS}{清廷围绕「白莲教战争持续扩展，…}嘉庆二年（1797年），清廷围绕「白莲教战争持续扩展，清军依赖乡勇、团练与地方资源」调遣军队、筹措饷械并处理战报，中央与地方的军政文书形成连续记录。核心取证：《清仁宗实录》嘉庆二年相关条；军机处档、战事方略及地方奏报。这一锚点记录的是文书链：它能证明呈报、上谕、部议或照会进入机构流程，不能由此推出实际执行。来源保留：此行仅确认相关文书链和可追查的行政动作；不得由其直接推断政策有效、地方服从、民众态度或单一因果。
\eventanchor{EQ-1798-MILITIA-MOBILIZATION-DOCUMENTS}{围绕「朝廷加强对乡勇、团练和地…}嘉庆三年（1798年），围绕「朝廷加强对乡勇、团练和地方绅士资源的动员以应对白莲教战争」的地方呈报、案件或赈务记录将社会反应带入官府文书链。核心取证：《清仁宗实录》嘉庆三年相关条；地方志、督抚奏折及地方档案。这一锚点记录的是文书链：它能证明呈报、上谕、部议或照会进入机构流程，不能由此推出实际执行。来源保留：此行仅确认相关文书链和可追查的行政动作；不得由其直接推断政策有效、地方服从、民众态度或单一因果。
\eventanchor{EQ-1799-HESHEN-PURGE-DOCUMENTS}{围绕「乾隆帝去世后，嘉庆帝逮治…}嘉庆四年（1799年），围绕「乾隆帝去世后，嘉庆帝逮治和珅并清理其政治网络」的上谕、奏折与部议在中央—地方行政链中流转。核心取证：《清仁宗实录》嘉庆四年相关条；宫中朱批奏折、内阁大库题本与《清史稿》相关志传。这一锚点记录的是文书链：它能证明呈报、上谕、部议或照会进入机构流程，不能由此推出实际执行。来源保留：此行仅确认相关文书链和可追查的行政动作；不得由其直接推断政策有效、地方服从、民众态度或单一因果。
\eventanchor{EQ-1797-WHITE-LOTUS-WAR-PRELUDE}{「白莲教战争持续扩展，清军依赖…}嘉庆二年（1797年），「白莲教战争持续扩展，清军依赖乡勇、团练与地方资源」之前的条件、信息和争议已通过中央与地方材料显现，构成该事件可追踪的前史节点。核心取证：《清仁宗实录》嘉庆二年相关条；军机处档、战事方略及地方奏报；同时期地方档案、地方志、日记或对方材料应按具体地区补检。这一锚点界定可回查的前史条件；条件、传闻或信息累积不应被写成已经证实的单一直接原因。来源保留：官修编年和地方材料的记录密度与立场并不相同；本行不将条件、传闻、呈报或后见解释直接等同为确定因果。
\eventanchor{EQ-1797-WHITE-LOTUS-WAR-DECISION}{围绕「白莲教战争持续扩展，清军…}嘉庆二年（1797年），围绕「白莲教战争持续扩展，清军依赖乡勇、团练与地方资源」的命令、调度、照会或呈报形成独立的中央—地方行动文书链。核心取证：《清仁宗实录》嘉庆二年相关条；军机处档、战事方略及地方奏报。这一锚点定位决策或调度动作；命令发出、地方接收、执行过程和实际结果仍须分别寻找证据。来源保留：奏折、上谕、部议、军机档或条约可证明行政动作和机构立场，不自动证明所有地区、群体或对手按其意图行动。
\eventanchor{EQ-1797-WHITE-LOTUS-WAR-AFTERMATH}{「白莲教战争持续扩展，清军依赖…}嘉庆二年（1797年），「白莲教战争持续扩展，清军依赖乡勇、团练与地方资源」引发的善后、执行或地方回应构成可由不同材料追踪的后续节点。核心取证：《清仁宗实录》嘉庆二年相关条；军机处档、战事方略及地方奏报；相关地区的档案、志书、契约、报刊或对方材料。这一锚点追踪后续影响；不同地区、群体与时段的反应不能由战后或改革后的概括性记忆代替。来源保留：战后、改革后或条约后的记忆、地方记录和官方总结常具有事后选择性；后续影响须按地点、群体和时间分层，不作整体化推断。
\eventanchor{EQ-1798-MILITIA-MOBILIZATION-PRELUDE}{「朝廷加强对乡勇、团练和地方绅…}嘉庆三年（1798年），「朝廷加强对乡勇、团练和地方绅士资源的动员以应对白莲教战争」之前的条件、信息和争议已通过中央与地方材料显现，构成该事件可追踪的前史节点。核心取证：《清仁宗实录》嘉庆三年相关条；地方志、督抚奏折及地方档案；同时期地方档案、地方志、日记或对方材料应按具体地区补检。这一锚点界定可回查的前史条件；条件、传闻或信息累积不应被写成已经证实的单一直接原因。来源保留：官修编年和地方材料的记录密度与立场并不相同；本行不将条件、传闻、呈报或后见解释直接等同为确定因果。
\subsection{1800—1809}
\eventanchor{EQ-1800-WHITE-LOTUS-FINANCE-DOCUMENTS}{围绕「白莲教战争军费、捐输与地…}嘉庆五年（1800年），围绕「白莲教战争军费、捐输与地方筹饷加重国家财政压力」的经费、税收、赈济或工程呈报经由户部与地方衙门处理。核心取证：《清仁宗实录》嘉庆五年相关条；户部奏销、盐政、河工或海关档案。这一锚点记录的是文书链：它能证明呈报、上谕、部议或照会进入机构流程，不能由此推出实际执行。来源保留：此行仅确认相关文书链和可追查的行政动作；不得由其直接推断政策有效、地方服从、民众态度或单一因果。
\eventanchor{EQ-1801-WHITE-LOTUS-SUPPRESSION-DOCUMENTS}{清廷围绕「清军在若干战区逐步压…}嘉庆六年（1801年），清廷围绕「清军在若干战区逐步压缩白莲教起事者活动空间」调遣军队、筹措饷械并处理战报，中央与地方的军政文书形成连续记录。核心取证：《清仁宗实录》嘉庆六年相关条；军机处档、战事方略及地方奏报。这一锚点记录的是文书链：它能证明呈报、上谕、部议或照会进入机构流程，不能由此推出实际执行。来源保留：此行仅确认相关文书链和可追查的行政动作；不得由其直接推断政策有效、地方服从、民众态度或单一因果。
\eventanchor{EQ-1802-WAR-REFUGEE-RELIEF-DOCUMENTS}{围绕「战区流民、复垦和赈济成为…}嘉庆七年（1802年），围绕「战区流民、复垦和赈济成为地方善后事务」的地方呈报、案件或赈务记录将社会反应带入官府文书链。核心取证：《清仁宗实录》嘉庆七年相关条；地方志、督抚奏折及地方档案。这一锚点记录的是文书链：它能证明呈报、上谕、部议或照会进入机构流程，不能由此推出实际执行。来源保留：此行仅确认相关文书链和可追查的行政动作；不得由其直接推断政策有效、地方服从、民众态度或单一因果。
\eventanchor{EQ-1803-FORBIDDEN-CITY-ATTACK-DOCUMENTS}{围绕「陈德等闯入紫禁城，暴露京…}嘉庆八年（1803年），围绕「陈德等闯入紫禁城，暴露京师警卫与内廷管理问题」的上谕、奏折与部议在中央—地方行政链中流转。核心取证：《清仁宗实录》嘉庆八年相关条；宫中朱批奏折、内阁大库题本与《清史稿》相关志传。这一锚点记录的是文书链：它能证明呈报、上谕、部议或照会进入机构流程，不能由此推出实际执行。来源保留：此行仅确认相关文书链和可追查的行政动作；不得由其直接推断政策有效、地方服从、民众态度或单一因果。
\eventanchor{EQ-1804-COASTAL-PIRACY-DOCUMENTS}{清廷围绕「东南海盗活动加剧，福…}嘉庆九年（1804年），清廷围绕「东南海盗活动加剧，福建、广东沿海军政压力上升」调遣军队、筹措饷械并处理战报，中央与地方的军政文书形成连续记录。核心取证：《清仁宗实录》嘉庆九年相关条；军机处档、战事方略及地方奏报。这一锚点记录的是文书链：它能证明呈报、上谕、部议或照会进入机构流程，不能由此推出实际执行。来源保留：此行仅确认相关文书链和可追查的行政动作；不得由其直接推断政策有效、地方服从、民众态度或单一因果。
\eventanchor{EQ-1805-WHITE-LOTUS-WAR-END-DOCUMENTS}{清廷围绕「白莲教大规模战争基本…}嘉庆十年（1805年），清廷围绕「白莲教大规模战争基本结束，清廷转入善后和治安重整」调遣军队、筹措饷械并处理战报，中央与地方的军政文书形成连续记录。核心取证：《清仁宗实录》嘉庆十年相关条；军机处档、战事方略及地方奏报。这一锚点记录的是文书链：它能证明呈报、上谕、部议或照会进入机构流程，不能由此推出实际执行。来源保留：此行仅确认相关文书链和可追查的行政动作；不得由其直接推断政策有效、地方服从、民众态度或单一因果。
\eventanchor{EQ-1806-TSAI-QIAN-PIRACY-DOCUMENTS}{清廷围绕「蔡牵等海盗集团持续活…}嘉庆十一年（1806年），清廷围绕「蔡牵等海盗集团持续活动，清廷调整水师与沿海防务」调遣军队、筹措饷械并处理战报，中央与地方的军政文书形成连续记录。核心取证：《清仁宗实录》嘉庆十一年相关条；军机处档、战事方略及地方奏报。这一锚点记录的是文书链：它能证明呈报、上谕、部议或照会进入机构流程，不能由此推出实际执行。来源保留：此行仅确认相关文书链和可追查的行政动作；不得由其直接推断政策有效、地方服从、民众态度或单一因果。
\eventanchor{EQ-1807-COASTAL-NAVAL-REFORM-DOCUMENTS}{围绕「清廷继续通过水师调度、招…}嘉庆十二年（1807年），围绕「清廷继续通过水师调度、招抚和地方协作应对海盗」的上谕、奏折与部议在中央—地方行政链中流转。核心取证：《清仁宗实录》嘉庆十二年相关条；宫中朱批奏折、内阁大库题本与《清史稿》相关志传。这一锚点记录的是文书链：它能证明呈报、上谕、部议或照会进入机构流程，不能由此推出实际执行。来源保留：此行仅确认相关文书链和可追查的行政动作；不得由其直接推断政策有效、地方服从、民众态度或单一因果。
\eventanchor{EQ-1808-PIRATE-ALLIANCES-DOCUMENTS}{围绕「郑一嫂、张保仔等海盗网络…}嘉庆十三年（1808年），围绕「郑一嫂、张保仔等海盗网络在华南海域形成强大联盟」的地方呈报、案件或赈务记录将社会反应带入官府文书链。核心取证：《清仁宗实录》嘉庆十三年相关条；地方志、督抚奏折及地方档案。这一锚点记录的是文书链：它能证明呈报、上谕、部议或照会进入机构流程，不能由此推出实际执行。来源保留：此行仅确认相关文书链和可追查的行政动作；不得由其直接推断政策有效、地方服从、民众态度或单一因果。
\eventanchor{EQ-1809-SOUTH-CHINA-SEA-CAMPAIGN-DOCUMENTS}{清廷围绕「清军、地方势力与外来…}嘉庆十四年（1809年），清廷围绕「清军、地方势力与外来船只在华南海域共同影响剿盗行动」调遣军队、筹措饷械并处理战报，中央与地方的军政文书形成连续记录。核心取证：《清仁宗实录》嘉庆十四年相关条；军机处档、战事方略及地方奏报。这一锚点记录的是文书链：它能证明呈报、上谕、部议或照会进入机构流程，不能由此推出实际执行。来源保留：此行仅确认相关文书链和可追查的行政动作；不得由其直接推断政策有效、地方服从、民众态度或单一因果。
\subsection{1810—1819}
\eventanchor{EQ-1810-ZHANG-BAOZAI-SURRENDER-DOCUMENTS}{清廷围绕「张保仔接受招安，华南…}嘉庆十五年（1810年），清廷围绕「张保仔接受招安，华南大海盗联盟瓦解」调遣军队、筹措饷械并处理战报，中央与地方的军政文书形成连续记录。核心取证：《清仁宗实录》嘉庆十五年相关条；军机处档、战事方略及地方奏报。这一锚点记录的是文书链：它能证明呈报、上谕、部议或照会进入机构流程，不能由此推出实际执行。来源保留：此行仅确认相关文书链和可追查的行政动作；不得由其直接推断政策有效、地方服从、民众态度或单一因果。
\eventanchor{EQ-1811-BANNER-LIVELIHOOD-DOCUMENTS}{围绕「八旗生计、债务和人口供养…}嘉庆十六年（1811年），围绕「八旗生计、债务和人口供养问题持续成为中央与地方的治理议题」的地方呈报、案件或赈务记录将社会反应带入官府文书链。核心取证：《清仁宗实录》嘉庆十六年相关条；地方志、督抚奏折及地方档案。这一锚点记录的是文书链：它能证明呈报、上谕、部议或照会进入机构流程，不能由此推出实际执行。来源保留：此行仅确认相关文书链和可追查的行政动作；不得由其直接推断政策有效、地方服从、民众态度或单一因果。
\eventanchor{EQ-1812-TIANLI-PREPARATION-DOCUMENTS}{围绕「天理教组织活动在直隶、河…}嘉庆十七年（1812年），围绕「天理教组织活动在直隶、河南、山东等地被官府逐步察觉」的地方呈报、案件或赈务记录将社会反应带入官府文书链。核心取证：《清仁宗实录》嘉庆十七年相关条；地方志、督抚奏折及地方档案。这一锚点记录的是文书链：它能证明呈报、上谕、部议或照会进入机构流程，不能由此推出实际执行。来源保留：此行仅确认相关文书链和可追查的行政动作；不得由其直接推断政策有效、地方服从、民众态度或单一因果。
\eventanchor{EQ-1813-TIANLI-UPRISING-DOCUMENTS}{清廷围绕「天理教起事者攻入紫禁…}嘉庆十八年（1813年），清廷围绕「天理教起事者攻入紫禁城部分区域，京师发生严重安全危机」调遣军队、筹措饷械并处理战报，中央与地方的军政文书形成连续记录。核心取证：《清仁宗实录》嘉庆十八年相关条；军机处档、战事方略及地方奏报。这一锚点记录的是文书链：它能证明呈报、上谕、部议或照会进入机构流程，不能由此推出实际执行。来源保留：此行仅确认相关文书链和可追查的行政动作；不得由其直接推断政策有效、地方服从、民众态度或单一因果。
\eventanchor{EQ-1814-TIANLI-SUPPRESSION-DOCUMENTS}{围绕「天理教起事被镇压，相关人…}嘉庆十九年（1814年），围绕「天理教起事被镇压，相关人员遭逮捕和处置」的上谕、奏折与部议在中央—地方行政链中流转。核心取证：《清仁宗实录》嘉庆十九年相关条；宫中朱批奏折、内阁大库题本与《清史稿》相关志传。这一锚点记录的是文书链：它能证明呈报、上谕、部议或照会进入机构流程，不能由此推出实际执行。来源保留：此行仅确认相关文书链和可追查的行政动作；不得由其直接推断政策有效、地方服从、民众态度或单一因果。
\eventanchor{EQ-1815-GRAND-CANAL-REPAIRS-DOCUMENTS}{围绕「漕运、河工和仓储问题在战…}嘉庆二十年（1815年），围绕「漕运、河工和仓储问题在战后财政压力下持续受到讨论」的经费、税收、赈济或工程呈报经由户部与地方衙门处理。核心取证：《清仁宗实录》嘉庆二十年相关条；户部奏销、盐政、河工或海关档案。这一锚点记录的是文书链：它能证明呈报、上谕、部议或照会进入机构流程，不能由此推出实际执行。来源保留：此行仅确认相关文书链和可追查的行政动作；不得由其直接推断政策有效、地方服从、民众态度或单一因果。
\eventanchor{EQ-1816-AMHERST-MISSION-DOCUMENTS}{围绕「阿美士德使团来华，礼仪与…}嘉庆二十一年（1816年），围绕「阿美士德使团来华，礼仪与外交交涉未获突破」的照会、条约文本、换约或地方执行报告分别进入外交文书链。核心取证：《清仁宗实录》嘉庆二十一年相关条；《筹办夷务始末》相关年卷与中外照会、条约文本。这一锚点记录的是文书链：它能证明呈报、上谕、部议或照会进入机构流程，不能由此推出实际执行。来源保留：此行仅确认相关文书链和可追查的行政动作；不得由其直接推断政策有效、地方服从、民众态度或单一因果。
\eventanchor{EQ-1817-SALT-ARREARS-DOCUMENTS}{围绕「盐课、商欠和财政亏空继续…}嘉庆二十二年（1817年），围绕「盐课、商欠和财政亏空继续困扰中央财政」的经费、税收、赈济或工程呈报经由户部与地方衙门处理。核心取证：《清仁宗实录》嘉庆二十二年相关条；户部奏销、盐政、河工或海关档案。这一锚点记录的是文书链：它能证明呈报、上谕、部议或照会进入机构流程，不能由此推出实际执行。来源保留：此行仅确认相关文书链和可追查的行政动作；不得由其直接推断政策有效、地方服从、民众态度或单一因果。
\eventanchor{EQ-1818-FLOOD-RELIEF-DOCUMENTS}{围绕「黄淮水患和赈务成为地方奏…}嘉庆二十三年（1818年），围绕「黄淮水患和赈务成为地方奏报与中央拨款的重要议题」的地方呈报、案件或赈务记录将社会反应带入官府文书链。核心取证：《清仁宗实录》嘉庆二十三年相关条；地方志、督抚奏折及地方档案。这一锚点记录的是文书链：它能证明呈报、上谕、部议或照会进入机构流程，不能由此推出实际执行。来源保留：此行仅确认相关文书链和可追查的行政动作；不得由其直接推断政策有效、地方服从、民众态度或单一因果。
\eventanchor{EQ-1819-OPIUM-TRADE-GROWTH-DOCUMENTS}{围绕「鸦片贸易增长引起官员对白…}嘉庆二十四年（1819年），围绕「鸦片贸易增长引起官员对白银流动、海防和社会秩序的担忧」的经费、税收、赈济或工程呈报经由户部与地方衙门处理。核心取证：《清仁宗实录》嘉庆二十四年相关条；户部奏销、盐政、河工或海关档案。这一锚点记录的是文书链：它能证明呈报、上谕、部议或照会进入机构流程，不能由此推出实际执行。来源保留：此行仅确认相关文书链和可追查的行政动作；不得由其直接推断政策有效、地方服从、民众态度或单一因果。
\subsection{1820—1829}
\eventanchor{EQ-1820-DAO-GUANG-ACCESSION-DOCUMENTS}{围绕「道光帝即位，财政、河工、…}嘉庆二十五年（1820年），围绕「道光帝即位，财政、河工、旗务和沿海贸易问题进入新朝」的上谕、奏折与部议在中央—地方行政链中流转。核心取证：《清仁宗实录》嘉庆二十五年相关条；宫中朱批奏折、内阁大库题本与《清史稿》相关志传。这一锚点记录的是文书链：它能证明呈报、上谕、部议或照会进入机构流程，不能由此推出实际执行。来源保留：此行仅确认相关文书链和可追查的行政动作；不得由其直接推断政策有效、地方服从、民众态度或单一因果。
\eventanchor{EQ-1821-IMPERIAL-ECONOMY-DOCUMENTS}{围绕「道光初年继续检讨内务府、…}道光元年（1821年），围绕「道光初年继续检讨内务府、官员经费和财政节用」的经费、税收、赈济或工程呈报经由户部与地方衙门处理。核心取证：《清宣宗实录》道光元年相关条；户部奏销、盐政、河工或海关档案。这一锚点记录的是文书链：它能证明呈报、上谕、部议或照会进入机构流程，不能由此推出实际执行。来源保留：此行仅确认相关文书链和可追查的行政动作；不得由其直接推断政策有效、地方服从、民众态度或单一因果。
\eventanchor{EQ-1822-YELLOW-RIVER-FLOOD-DOCUMENTS}{围绕「黄河水患、堤工和灾赈问题…}道光二年（1822年），围绕「黄河水患、堤工和灾赈问题反复进入中央议程」的地方呈报、案件或赈务记录将社会反应带入官府文书链。核心取证：《清宣宗实录》道光二年相关条；地方志、督抚奏折及地方档案。这一锚点记录的是文书链：它能证明呈报、上谕、部议或照会进入机构流程，不能由此推出实际执行。来源保留：此行仅确认相关文书链和可追查的行政动作；不得由其直接推断政策有效、地方服从、民众态度或单一因果。
\eventanchor{EQ-1823-BANNER-DEBT-DOCUMENTS}{围绕「旗人债务、典当和生计问题…}道光三年（1823年），围绕「旗人债务、典当和生计问题持续引起官府讨论」的地方呈报、案件或赈务记录将社会反应带入官府文书链。核心取证：《清宣宗实录》道光三年相关条；地方志、督抚奏折及地方档案。这一锚点记录的是文书链：它能证明呈报、上谕、部议或照会进入机构流程，不能由此推出实际执行。来源保留：此行仅确认相关文书链和可追查的行政动作；不得由其直接推断政策有效、地方服从、民众态度或单一因果。
\eventanchor{EQ-1824-GRAND-CANAL-CRISIS-DOCUMENTS}{围绕「运河水运和漕粮转输面临河…}道光四年（1824年），围绕「运河水运和漕粮转输面临河工、淤积与财政问题」的经费、税收、赈济或工程呈报经由户部与地方衙门处理。核心取证：《清宣宗实录》道光四年相关条；户部奏销、盐政、河工或海关档案。这一锚点记录的是文书链：它能证明呈报、上谕、部议或照会进入机构流程，不能由此推出实际执行。来源保留：此行仅确认相关文书链和可追查的行政动作；不得由其直接推断政策有效、地方服从、民众态度或单一因果。
\eventanchor{EQ-1825-SOUTHWEST-MINING-UNREST-DOCUMENTS}{围绕「云贵矿业、移民和治安纠纷…}道光五年（1825年），围绕「云贵矿业、移民和治安纠纷继续考验地方治理」的地方呈报、案件或赈务记录将社会反应带入官府文书链。核心取证：《清宣宗实录》道光五年相关条；地方志、督抚奏折及地方档案。这一锚点记录的是文书链：它能证明呈报、上谕、部议或照会进入机构流程，不能由此推出实际执行。来源保留：此行仅确认相关文书链和可追查的行政动作；不得由其直接推断政策有效、地方服从、民众态度或单一因果。
\eventanchor{EQ-1826-JAHANGIR-KHOJA-UPRISING-DOCUMENTS}{清廷围绕「张格尔和卓进入喀什噶…}道光六年（1826年），清廷围绕「张格尔和卓进入喀什噶尔，天山南路爆发大规模反清动荡」调遣军队、筹措饷械并处理战报，中央与地方的军政文书形成连续记录。核心取证：《清宣宗实录》道光六年相关条；军机处档、战事方略及地方奏报。这一锚点记录的是文书链：它能证明呈报、上谕、部议或照会进入机构流程，不能由此推出实际执行。来源保留：此行仅确认相关文书链和可追查的行政动作；不得由其直接推断政策有效、地方服从、民众态度或单一因果。
\eventanchor{EQ-1827-KASHGAR-RECOVERY-DOCUMENTS}{清廷围绕「清军反攻并收复喀什噶…}道光七年（1827年），清廷围绕「清军反攻并收复喀什噶尔等地，张格尔势力受挫」调遣军队、筹措饷械并处理战报，中央与地方的军政文书形成连续记录。核心取证：《清宣宗实录》道光七年相关条；军机处档、战事方略及地方奏报。这一锚点记录的是文书链：它能证明呈报、上谕、部议或照会进入机构流程，不能由此推出实际执行。来源保留：此行仅确认相关文书链和可追查的行政动作；不得由其直接推断政策有效、地方服从、民众态度或单一因果。
\eventanchor{EQ-1828-JAHANGIR-CAPTURE-DOCUMENTS}{清廷围绕「张格尔被捕并押解处置…}道光八年（1828年），清廷围绕「张格尔被捕并押解处置，清廷加强南疆军政控制」处理盟旗、驻防、交通或地方呈报，相关边务文书进入中央决策链。核心取证：《清宣宗实录》道光八年相关条；军机处档、理藩院档或相关方略。这一锚点记录的是文书链：它能证明呈报、上谕、部议或照会进入机构流程，不能由此推出实际执行。来源保留：此行仅确认相关文书链和可追查的行政动作；不得由其直接推断政策有效、地方服从、民众态度或单一因果。
\eventanchor{EQ-1829-XINJIANG-FINANCE-DOCUMENTS}{围绕「南疆战后军费、驻防与贸易…}道光九年（1829年），围绕「南疆战后军费、驻防与贸易秩序成为清廷边务财政负担」的经费、税收、赈济或工程呈报经由户部与地方衙门处理。核心取证：《清宣宗实录》道光九年相关条；户部奏销、盐政、河工或海关档案。这一锚点记录的是文书链：它能证明呈报、上谕、部议或照会进入机构流程，不能由此推出实际执行。来源保留：此行仅确认相关文书链和可追查的行政动作；不得由其直接推断政策有效、地方服从、民众态度或单一因果。
\subsection{1830—1839}
\eventanchor{EQ-1830-RUSSO-QING-TRADE-DOCUMENTS}{围绕「恰克图等边境贸易和俄清交…}道光十年（1830年），围绕「恰克图等边境贸易和俄清交涉持续影响北部边务」的照会、条约文本、换约或地方执行报告分别进入外交文书链。核心取证：《清宣宗实录》道光十年相关条；《筹办夷务始末》相关年卷与中外照会、条约文本。这一锚点记录的是文书链：它能证明呈报、上谕、部议或照会进入机构流程，不能由此推出实际执行。来源保留：此行仅确认相关文书链和可追查的行政动作；不得由其直接推断政策有效、地方服从、民众态度或单一因果。
\eventanchor{EQ-1831-CENTRAL-ASIA-CHOLERA-DOCUMENTS}{围绕「疾病流行与交通网络对军队…}道光十一年（1831年），围绕「疾病流行与交通网络对军队、城市和商路造成压力」的地方呈报、案件或赈务记录将社会反应带入官府文书链。核心取证：《清宣宗实录》道光十一年相关条；地方志、督抚奏折及地方档案。这一锚点记录的是文书链：它能证明呈报、上谕、部议或照会进入机构流程，不能由此推出实际执行。来源保留：此行仅确认相关文书链和可追查的行政动作；不得由其直接推断政策有效、地方服从、民众态度或单一因果。
\eventanchor{EQ-1832-COASTAL-DEFENCE-DOCUMENTS}{围绕「鸦片走私增长背景下的海防…}道光十二年（1832年），围绕「鸦片走私增长背景下的海防、缉私和地方官责任成为争议」的上谕、奏折与部议在中央—地方行政链中流转。核心取证：《清宣宗实录》道光十二年相关条；宫中朱批奏折、内阁大库题本与《清史稿》相关志传。这一锚点记录的是文书链：它能证明呈报、上谕、部议或照会进入机构流程，不能由此推出实际执行。来源保留：此行仅确认相关文书链和可追查的行政动作；不得由其直接推断政策有效、地方服从、民众态度或单一因果。
\eventanchor{EQ-1833-OPIUM-POLICY-DEBATE-DOCUMENTS}{围绕「官员围绕禁烟、征税、白银…}道光十三年（1833年），围绕「官员围绕禁烟、征税、白银和海防提出相互竞争的方案」的上谕、奏折与部议在中央—地方行政链中流转。核心取证：《清宣宗实录》道光十三年相关条；宫中朱批奏折、内阁大库题本与《清史稿》相关志传。这一锚点记录的是文书链：它能证明呈报、上谕、部议或照会进入机构流程，不能由此推出实际执行。来源保留：此行仅确认相关文书链和可追查的行政动作；不得由其直接推断政策有效、地方服从、民众态度或单一因果。
\eventanchor{EQ-1834-NAPIER-AFFAIR-DOCUMENTS}{围绕「律劳卑抵广州后与两广官员…}道光十四年（1834年），围绕「律劳卑抵广州后与两广官员发生外交摩擦」的照会、条约文本、换约或地方执行报告分别进入外交文书链。核心取证：《清宣宗实录》道光十四年相关条；《筹办夷务始末》相关年卷与中外照会、条约文本。这一锚点记录的是文书链：它能证明呈报、上谕、部议或照会进入机构流程，不能由此推出实际执行。来源保留：此行仅确认相关文书链和可追查的行政动作；不得由其直接推断政策有效、地方服从、民众态度或单一因果。
\eventanchor{EQ-1835-GUANGZHOU-TRADE-REGULATION-DOCUMENTS}{围绕「广州贸易、行商债务和外商…}道光十五年（1835年），围绕「广州贸易、行商债务和外商管理继续在地方与中央之间协商」的经费、税收、赈济或工程呈报经由户部与地方衙门处理。核心取证：《清宣宗实录》道光十五年相关条；户部奏销、盐政、河工或海关档案。这一锚点记录的是文书链：它能证明呈报、上谕、部议或照会进入机构流程，不能由此推出实际执行。来源保留：此行仅确认相关文书链和可追查的行政动作；不得由其直接推断政策有效、地方服从、民众态度或单一因果。
\eventanchor{EQ-1836-OPIUM-LEGALIZATION-DEBATE-DOCUMENTS}{围绕「许乃济等围绕鸦片弛禁、征…}道光十六年（1836年），围绕「许乃济等围绕鸦片弛禁、征税与禁绝提出政策论争」的上谕、奏折与部议在中央—地方行政链中流转。核心取证：《清宣宗实录》道光十六年相关条；宫中朱批奏折、内阁大库题本与《清史稿》相关志传。这一锚点记录的是文书链：它能证明呈报、上谕、部议或照会进入机构流程，不能由此推出实际执行。来源保留：此行仅确认相关文书链和可追查的行政动作；不得由其直接推断政策有效、地方服从、民众态度或单一因果。
\eventanchor{EQ-1837-CANTON-SMUGGLING-DOCUMENTS}{围绕「广东缉私与鸦片走私之间的…}道光十七年（1837年），围绕「广东缉私与鸦片走私之间的冲突持续升级」的地方呈报、案件或赈务记录将社会反应带入官府文书链。核心取证：《清宣宗实录》道光十七年相关条；地方志、督抚奏折及地方档案。这一锚点记录的是文书链：它能证明呈报、上谕、部议或照会进入机构流程，不能由此推出实际执行。来源保留：此行仅确认相关文书链和可追查的行政动作；不得由其直接推断政策有效、地方服从、民众态度或单一因果。
\eventanchor{EQ-1838-LIN-ZEXU-APPOINTMENT-DOCUMENTS}{围绕「林则徐受命赴广东查禁鸦片…}道光十八年（1838年），围绕「林则徐受命赴广东查禁鸦片，禁烟政策转为高压执行」的上谕、奏折与部议在中央—地方行政链中流转。核心取证：《清宣宗实录》道光十八年相关条；宫中朱批奏折、内阁大库题本与《清史稿》相关志传。这一锚点记录的是文书链：它能证明呈报、上谕、部议或照会进入机构流程，不能由此推出实际执行。来源保留：此行仅确认相关文书链和可追查的行政动作；不得由其直接推断政策有效、地方服从、民众态度或单一因果。
\eventanchor{EQ-1839-HUMEN-OPIUM-DESTRUCTION-DOCUMENTS}{围绕「林则徐在虎门销毁收缴鸦片…}道光十九年（1839年），围绕「林则徐在虎门销毁收缴鸦片，禁烟冲突公开化」的照会、条约文本、换约或地方执行报告分别进入外交文书链。核心取证：《清宣宗实录》道光十九年相关条；《筹办夷务始末》相关年卷与中外照会、条约文本。这一锚点记录的是文书链：它能证明呈报、上谕、部议或照会进入机构流程，不能由此推出实际执行。来源保留：此行仅确认相关文书链和可追查的行政动作；不得由其直接推断政策有效、地方服从、民众态度或单一因果。
\subsection{1840—1849}
\eventanchor{EQ-1840-FIRST-OPIUM-WAR-DOCUMENTS}{清廷围绕「英国舰队来华，第一次…}道光二十年（1840年），清廷围绕「英国舰队来华，第一次鸦片战争爆发」调遣军队、筹措饷械并处理战报，中央与地方的军政文书形成连续记录。核心取证：《清宣宗实录》道光二十年相关条；军机处档、战事方略及地方奏报。这一锚点记录的是文书链：它能证明呈报、上谕、部议或照会进入机构流程，不能由此推出实际执行。来源保留：此行仅确认相关文书链和可追查的行政动作；不得由其直接推断政策有效、地方服从、民众态度或单一因果。
\eventanchor{EQ-1841-CHUANBI-NEGOTIATIONS-DOCUMENTS}{围绕「穿鼻草约等谈判和沿海战事…}道光二十一年（1841年），围绕「穿鼻草约等谈判和沿海战事并行，清廷内部对和战意见分歧」的照会、条约文本、换约或地方执行报告分别进入外交文书链。核心取证：《清宣宗实录》道光二十一年相关条；《筹办夷务始末》相关年卷与中外照会、条约文本。这一锚点记录的是文书链：它能证明呈报、上谕、部议或照会进入机构流程，不能由此推出实际执行。来源保留：此行仅确认相关文书链和可追查的行政动作；不得由其直接推断政策有效、地方服从、民众态度或单一因果。
\eventanchor{EQ-1841-CANTON-BATTLES-DOCUMENTS}{清廷围绕「广州周边战事和赔款谈…}道光二十一年（1841年），清廷围绕「广州周边战事和赔款谈判反复进行」调遣军队、筹措饷械并处理战报，中央与地方的军政文书形成连续记录。核心取证：《清宣宗实录》道光二十一年相关条；军机处档、战事方略及地方奏报。这一锚点记录的是文书链：它能证明呈报、上谕、部议或照会进入机构流程，不能由此推出实际执行。来源保留：此行仅确认相关文书链和可追查的行政动作；不得由其直接推断政策有效、地方服从、民众态度或单一因果。
\eventanchor{EQ-1842-NANJING-TREATY-DOCUMENTS}{围绕「《南京条约》签订，清廷割…}道光二十二年（1842年），围绕「《南京条约》签订，清廷割让香港并开放五口通商、支付赔款」的照会、条约文本、换约或地方执行报告分别进入外交文书链。核心取证：《清宣宗实录》道光二十二年相关条；《筹办夷务始末》相关年卷与中外照会、条约文本。这一锚点记录的是文书链：它能证明呈报、上谕、部议或照会进入机构流程，不能由此推出实际执行。来源保留：此行仅确认相关文书链和可追查的行政动作；不得由其直接推断政策有效、地方服从、民众态度或单一因果。
\eventanchor{EQ-1843-BOGUE-TREATY-DOCUMENTS}{围绕「《虎门条约》补充领事裁判…}道光二十三年（1843年），围绕「《虎门条约》补充领事裁判、协定关税等安排」的照会、条约文本、换约或地方执行报告分别进入外交文书链。核心取证：《清宣宗实录》道光二十三年相关条；《筹办夷务始末》相关年卷与中外照会、条约文本。这一锚点记录的是文书链：它能证明呈报、上谕、部议或照会进入机构流程，不能由此推出实际执行。来源保留：此行仅确认相关文书链和可追查的行政动作；不得由其直接推断政策有效、地方服从、民众态度或单一因果。
\eventanchor{EQ-1844-WANGXIA-WHAMPOA-TREATIES-DOCUMENTS}{围绕「《望厦条约》《黄埔条约》…}道光二十四年（1844年），围绕「《望厦条约》《黄埔条约》分别确立美法在华通商与相关权利」的照会、条约文本、换约或地方执行报告分别进入外交文书链。核心取证：《清宣宗实录》道光二十四年相关条；《筹办夷务始末》相关年卷与中外照会、条约文本。这一锚点记录的是文书链：它能证明呈报、上谕、部议或照会进入机构流程，不能由此推出实际执行。来源保留：此行仅确认相关文书链和可追查的行政动作；不得由其直接推断政策有效、地方服从、民众态度或单一因果。
\eventanchor{EQ-1845-TREATY-PORT-OPENING-DOCUMENTS}{围绕「五口通商制度进入具体开埠…}道光二十五年（1845年），围绕「五口通商制度进入具体开埠、税则和城市管理阶段」的经费、税收、赈济或工程呈报经由户部与地方衙门处理。核心取证：《清宣宗实录》道光二十五年相关条；户部奏销、盐政、河工或海关档案。这一锚点记录的是文书链：它能证明呈报、上谕、部议或照会进入机构流程，不能由此推出实际执行。来源保留：此行仅确认相关文书链和可追查的行政动作；不得由其直接推断政策有效、地方服从、民众态度或单一因果。
\eventanchor{EQ-1846-MISSIONARY-RETURN-DOCUMENTS}{围绕「条约环境下传教活动与地方…}道光二十六年（1846年），围绕「条约环境下传教活动与地方教案问题逐渐增多」的地方呈报、案件或赈务记录将社会反应带入官府文书链。核心取证：《清宣宗实录》道光二十六年相关条；地方志、督抚奏折及地方档案。这一锚点记录的是文书链：它能证明呈报、上谕、部议或照会进入机构流程，不能由此推出实际执行。来源保留：此行仅确认相关文书链和可追查的行政动作；不得由其直接推断政策有效、地方服从、民众态度或单一因果。
\eventanchor{EQ-1847-GUANGZHOU-CONFLICT-DOCUMENTS}{围绕「广州周边的中英冲突显示条…}道光二十七年（1847年），围绕「广州周边的中英冲突显示条约执行和城市主权问题尚未解决」的照会、条约文本、换约或地方执行报告分别进入外交文书链。核心取证：《清宣宗实录》道光二十七年相关条；《筹办夷务始末》相关年卷与中外照会、条约文本。这一锚点记录的是文书链：它能证明呈报、上谕、部议或照会进入机构流程，不能由此推出实际执行。来源保留：此行仅确认相关文书链和可追查的行政动作；不得由其直接推断政策有效、地方服从、民众态度或单一因果。
\eventanchor{EQ-1848-YANGTZE-TRADE-DOCUMENTS}{围绕「通商口岸贸易网络向长江和…}道光二十八年（1848年），围绕「通商口岸贸易网络向长江和沿海延伸，地方商人开始适应新制度」的经费、税收、赈济或工程呈报经由户部与地方衙门处理。核心取证：《清宣宗实录》道光二十八年相关条；户部奏销、盐政、河工或海关档案。这一锚点记录的是文书链：它能证明呈报、上谕、部议或照会进入机构流程，不能由此推出实际执行。来源保留：此行仅确认相关文书链和可追查的行政动作；不得由其直接推断政策有效、地方服从、民众态度或单一因果。
\eventanchor{EQ-1849-GUANGZHOU-CITY-ENTRY-DOCUMENTS}{围绕「英国要求进入广州城的争议…}道光二十九年（1849年），围绕「英国要求进入广州城的争议再次激化，地方民众与官府关系受影响」的地方呈报、案件或赈务记录将社会反应带入官府文书链。核心取证：《清宣宗实录》道光二十九年相关条；地方志、督抚奏折及地方档案。这一锚点记录的是文书链：它能证明呈报、上谕、部议或照会进入机构流程，不能由此推出实际执行。来源保留：此行仅确认相关文书链和可追查的行政动作；不得由其直接推断政策有效、地方服从、民众态度或单一因果。
\subsection{1850—1859}
\eventanchor{EQ-1850-XIANFENG-ACCESSION-DOCUMENTS}{围绕「咸丰帝即位，财政紧张、内…}道光三十年（1850年），围绕「咸丰帝即位，财政紧张、内乱和对外压力同时加重」的上谕、奏折与部议在中央—地方行政链中流转。核心取证：《清宣宗实录》道光三十年相关条；宫中朱批奏折、内阁大库题本与《清史稿》相关志传。这一锚点记录的是文书链：它能证明呈报、上谕、部议或照会进入机构流程，不能由此推出实际执行。来源保留：此行仅确认相关文书链和可追查的行政动作；不得由其直接推断政策有效、地方服从、民众态度或单一因果。
